\section{Compact measure constructions used in the proof}
\label{sec:measure-constructions}

The spectral argument and the common signature model use compactness
only for compact metric spaces, countable arrays, and bounded finite
measures. We record constructions for these particular uses, including
the jointly measurable densities in \Cref{lem:limit-model}.

\begin{lemma}[Compact arrays and finite measures]
\label{lem:compact-measures}
Let $X$ be a compact metric space and $C<\infty$.
The positive Borel measures on $X$ of mass at most $C$ form a compact
metrizable space for convergence against continuous functions.
Countable products of compact metric spaces are compact metrizable.
If $\nu_j$ converges weakly to $\nu$ and $F\subseteq X$ is closed,
then $\limsup_j\nu_j(F)\le\nu(F)$.
\end{lemma}
\begin{proof}
Choose nested finite Borel partitions of $X$ whose cell diameters tend
to zero. Such partitions result by taking a finite cover by balls of
arbitrarily small radius, making it disjoint by successive differences,
and refining the previous partition. Given a sequence of measures,
successive subsequences make every cell mass converge, as well as the
total mass. The limiting cell masses are nonnegative and consistent
under refinement.

Realize these masses by subdividing an interval of the limiting total
length, first into intervals of the masses of the first partition,
then recursively into intervals for its refinements. Outside the
countable set of subdivision endpoints, each interval point specifies
a nested sequence of cells. Representatives from those cells form a
Cauchy sequence: all representatives after level $j$ belong to one
cell of diameter tending to zero. Compactness gives a limiting point
of $X$. The representative step maps are measurable, so their limit
is measurable. Push forward interval length by this limit map.
Uniform continuity shows that the integral of each continuous function
is the limit of its finite representative sums. The same sums
approximate the integrals of the original measures uniformly. This
proves weak subsequential compactness, including zero limiting mass.

Continuous functions on $X$ have a countable uniformly dense family.
For instance, finite nets from a countable dense subset of $X$, together
with continuous distance cutoffs and rational values, give uniformly
accurate piecewise weighted approximations to each continuous function.
If $(f_j)$ is such a family, convergence of the integrals of all $f_j$
metrizes weak convergence on measures of bounded mass, by uniform
approximation. The preceding subsequence argument therefore gives
compactness in that metric. For a countable product use the metric
$\sum_j2^{-j}\min\{1,d_j\}$; successive coordinate subsequences and a
diagonal choice prove its compactness.

Finally the continuous functions
$f_k(x)=\max\{0,1-k\,d(x,F)\}$ decrease to $\mathbf1_F$.
For every $k$, weak convergence gives
$\limsup_j\nu_j(F)\le\int f_k\,d\nu$.
Bounded monotone approximation as $k\to\infty$ gives the stated
closed-set inequality.
\end{proof}

In particular, consistent masses of finite bit cylinders define a
measure on the compact bit-sequence space: apply the interval
subdivision construction directly to the successive bit prefixes.
If the masses depend measurably on a stored record, all subdivision
thresholds are measurable functions of that record. Using an additional
uniform interval coordinate gives a jointly measurable sample. The same
construction works for a compact metric target using the nested
partitions above: evaluation of a measure on a Borel cell is measurable
in the weak topology. For closed cells this follows from continuous
distance cutoffs, and a monotone-class argument gives all Borel cells. Two
independent interval coordinates sample independently from the two
conditional measures stored in \Cref{lem:limit-model}. Thus the argument
retains those conditional measures before taking limits; it does not
interchange an unrecorded conditional distribution with a weak limit.

\begin{lemma}[Densities of dominated cylinder measures]
\label{lem:cylinder-densities}
Let $X$ be a compact bit-sequence space with Borel probability $m$,
let $\theta$ have a probability law, and suppose $\nu_\theta$ is a
measurable family of measures satisfying $0\le\nu_\theta\le m$.
There is a jointly measurable $h(\theta,x)\in[0,1]$ such that,
for almost every $\theta$, $\nu_\theta(A)=\int_Ah(\theta,x)\,dm(x)$
for every Borel $A$. Conditional averages of an $L^2(m)$ function
on increasing finite-prefix partitions converge in $L^2(m)$.
\end{lemma}
\begin{proof}
Let $\mathcal P_k$ be the partition by the first $k$ bits, and set
$h_k(\theta,x)=\nu_\theta(C)/m(C)$ when $x\in C\in\mathcal P_k$
and $m(C)>0$, with value zero on null cells. These functions are
jointly measurable and lie in $[0,1]$. Their successive differences
are orthogonal in $L^2(d\Pr(\theta)\,dm(x))$, because the average of
$h_{k+1}-h_k$ over each parent cell is zero for every $\theta$.
The squared norms of $h_k$ are bounded by one. Hence $(h_k)$ is
Cauchy in $L^2$, and has an $L^2$ limit $h$. One can see completeness
here directly: choose a subsequence with summable $L^2$ norms of its
successive differences; their $L^1$ norms are also summable, giving
an almost-everywhere convergent sum with the required $L^2$ tail bound.
A further almost-everywhere convergent subsequence makes $h$ jointly
measurable with values in $[0,1]$.

For each cylinder $C$, all sufficiently late $h_k$ integrate over $C$
to $\nu_\theta(C)$. Convergence in the product $L^2$ space gives
the same identity for $h$, first in $L^1$ of $\theta$ and thus almost
surely. There are countably many cylinders, so these identities hold
simultaneously outside one null set. They then extend to all Borel
sets: the class of sets approximable in measure by the finite cylinder
algebra is a sigma-algebra, and both measures are dominated by $m$.

Finite-cell averaging is an $L^2$ contraction by Cauchy--Schwarz.
Cylinder-simple functions are dense in $L^2(m)$, by simple-function
approximation and the same cylinder approximation of measurable sets.
Approximating a function by such a simple function and using the
contraction proves the final assertion.
\end{proof}

Rescaling gives the same density construction for $\nu_\theta\le Rm$.
For countably many ordered measures, the versions can be chosen
simultaneously with their almost-everywhere ordering; their increasing
supremum represents the increasing supremum of the measures. This is
the use in \Cref{lem:limit-zero-product}. The integration identity for
the sampled kernels holds first for rectangle indicators. The class of
sets satisfying it is closed under complements and disjoint countable
unions; hence it holds on the product sigma-algebra and then for
nonnegative or bounded measurable integrands by simple-function
approximation. These constructions justify the limiting integrations
and projections used in \Cref{sec:common-limit}.
